\emph{Set-up.} Take $x_{0} = 0$ by a translation of the labels. Write $x = ly$, $x' = ly'$, $d = y - y'$, $\bar y = \tfrac12(y + y')$, $0 < l \le l_{1}$, with $y, y'$ in a ball $B$ such that $lB \subset N$ for $l \le l_{1}$, $N$ a geodesically convex neighbourhood of $x_{0}$ for the derived metric $g_{\mu\nu} = g_{\mu}\cdot g_{\nu}$, $g_{\mu} = \hb^{-1}(e_{\mu})$ (Definition \claimref{obj.synthesis.gtg-flat-spacetime-gauge-fields}(a), (d); Fact \claimref{synthesis.imported-gtg-dictionary}). The Dirac field and its CAR algebra are those of Fact \claimref{synthesis.imported-dimock-dirac-car-algebra-on-a-gravitational-field}, with the flat-field conventions of Proposition \claimref{field-theory.second-quantisation}. The scaled field is $\psi^{l}(y) = l^{3/2}\psi(ly)$, so its two-point functions are $l^{3}$ times those of $\psi$ at $(ly, ly')$. Matrix and index forms below are readings of the spacetime-algebra statements.

\emph{Step 1 (the Hadamard form; imported and cited from memory, written in the spacetime algebra).} By Fact \claimref{synthesis.imported-dirac-hadamard-form-propagates-and-adiabatic-states-are-hadamard}(a) the two-point functions of a Hadamard state (Fact \claimref{synthesis.imported-microlocal-hadamard-condition}(b)) have the Hadamard form of Sahlmann and Verch (the Fact is their Thm.~5.8, which according to the introduction of \cite{SahlmannVerch2001} repairs a gap in Radzikowski's argument that also affects the Dirac equivalence proofs of Koehler, Kratzert and Hollands; the statement of Thm.~5.8 was not checked against the text; Sanders \cite{Sanders2010} also addresses gaps in the Dirac equivalence, cited from memory): the Dirac operator applied to a distribution of Hadamard form for the squared operator, plus a remainder that, for the truncation of order $n$, is of class $C^{n}$. The squared operator $(R_{I_{4}}D - m_{0})(R_{I_{4}}D + m_{0}) = -(D^{2} + m_{0}^{2})$ ($R_{I_{4}} : x \mapsto xI_{4}$ the field theory's unit acting on the right, commuting with every left factor; numbers are read in $\mathrm{span}\{1, I_{4}\}$ and act on the right in the same way) is normally hyperbolic, and by the Lichnerowicz formula its connection is the spin connection $D_{a}\psi = a\cdot\nabla\psi + \tfrac12\Omega(a)\psi$. For such an operator (standard, Fact \claimref{synthesis.imported-hadamard-parametrix-of-the-squared-dirac-operator-leads-with-van-vleck-root-times-spin-transport}(a), (b): Friedlander 1975; Bar, Ginoux and Pfaffle 2007, Ch.~2) the Hadamard parametrix on $N$ is $-\tfrac{1}{8\pi^{2}}(U/\sigma + V\log\sigma + W)$ with $\sigma$ the Synge function, $U = \Delta^{1/2}P$, $P$ the parallel transport of the connection along the geodesic from $x'$ to $x$, $V, W$ smooth (on the flat field $V(x, x) = -\tfrac12 m_{0}^{2}$, the value for which $(\Box + m_{0}^{2})(1/\sigma + V\log\sigma)$ has no $1/\sigma$ term, since $\Box\log\sigma = \Box\sigma/\sigma - Q(\nabla\sigma)/\sigma^{2} = 2/\sigma$; computed); the constant $-\tfrac{1}{8\pi^{2}}$ is fixed by the flat case (Step 6). Parallel transport of a spinor along $\gamma$ solves $\dot\gamma\cdot\nabla\psi + \tfrac12\Omega(\dot\gamma)\psi = 0$, so $P$ is left multiplication by the rotor $R$ with $\dot R = -\tfrac12\Omega(\dot\gamma)R$, $R = 1$ at $t = 0$ ($R\tilde R = 1$ since $\Omega(\dot\gamma)$ is a bivector). The boundary values $\sigma_{\pm}$ carry the $I_{4}0$ of a time function, an infinitesimal shift of the separation written in scalar form; the limit does not depend on which time function (Fact \claimref{synthesis.imported-hadamard-parametrix-of-the-squared-dirac-operator-leads-with-van-vleck-root-times-spin-transport}(c): Kay and Wald 1991; Sahlmann and Verch, Sec.~5.1). Positive frequency is fixed as in the stored flat kernels (Theorem \claimref{synthesis.bulk-hadamard-states-restrict-to-the-generator-condition-on-flat-null-planes}(ii)): $\omega(\psi\psi^{\dagger})$ carries $(p + m_{0})e_{0}e^{-I_{4}p\cdot z}$, $p = R_{I_{4}}\nabla$ on $e^{-I_{4}p\cdot z}$, whence $(R_{I_{4}}D_{x} + m_{0})H_{+}e_{0}$; the other ordering carries $(p - m_{0})e_{0}e^{+I_{4}p\cdot z}$, $p = -R_{I_{4}}\nabla$ there, whence $-(R_{I_{4}}D_{x} + m_{0})H_{-}e_{0}$. The fibre form is the Dirac adjoint, carried by $\psi^{\dagger}$ and the factor $e_{0}$ of the stored kernels (matrix reading $\bar\psi = \psi^{\dagger}\gamma_{0}$), the same at every point because Gauge Theory Gravity trivialises the spinors by the background. In Sahlmann and Verch the matrix $\Theta$ is that of the Dirac conjugation composed with charge conjugation, as the Majorana pairing needs; in the gauge-invariant Dirac form $\omega(\psi\psi^{\dagger})$ its role is taken by this constant $e_{0}$. Of the remainder only $C^{1}$ is used ($n = 1$). Gauge invariance enters here (argued, not checked): the source of the Fact is posed for Majorana spinors on the self-dual CAR algebra, and a Dirac field is a doubled Majorana field whose two-point function contains the anomalous blocks $\omega(\psi\psi)$, $\omega(\psi^{\dagger}\psi^{\dagger})$; for a gauge-invariant state these vanish and the doubled two-point function is block-diagonal with blocks the store's pair $\mathcal{G}^{(\pm)}$, so the Fact applies to them as stated.

\emph{Step 2 (the Synge function; proved, computed).} Let $\gamma(t) = l\eta(t)$, $t \in [0, 1]$, be the geodesic from $ly'$ to $ly$, the worldline whose velocity $\hb^{-1}(\dot\gamma)$ is $D$-parallel (Fact \claimref{synthesis.imported-gtg-dictionary}(iii)). In the scaled labels this reads $\ddot\eta = -l\,\Gamma(l\eta)(\dot\eta, \dot\eta)$, with $\Gamma$ the Christoffel symbols $g^{\rho}\cdot D_{e_{\mu}}g_{\nu}$ of the derived metric as its component reading (Fact \claimref{synthesis.imported-gtg-dictionary}(ii)), $\eta(0) = y'$, $\eta(1) = y$, whose solution is smooth in $(l, y, y')$ down to $l = 0$ with $\eta = y' + td + O(l)$ (for $l\Gamma$ small the boundary-value problem is a contraction). Hence $\sigma = \tfrac12\int_{0}^{1}g(\dot\gamma, \dot\gamma)dt = l^{2}\sigma_{l}$, $\sigma_{l} = \tfrac12\int_{0}^{1}g(l\eta)(\dot\eta, \dot\eta)dt$, jointly smooth, and
\[
\sigma_{l} = \tfrac12 Q(d) + l\,d\cdot\bigl[(\bar y\cdot\nabla)\hb^{-1}\bigr](d)\big|_{0} + O(l^{2}),
\]
(index reading $\tfrac{l}{2}(\partial_{\rho}g_{\mu\nu})(0)\bar y^{\rho}d^{\mu}d^{\nu}$, since $g(a, b) = \hb^{-1}(a)\cdot\hb^{-1}(b)$ gives $(\bar y\cdot\nabla)g(d, d) = 2\,\hb^{-1}(d)\cdot[(\bar y\cdot\nabla)\hb^{-1}](d)$), using $g(0) = \eta_{\mu\nu}$ ($\hb(0) = 1$), $\int_{0}^{1}(1 - 2t)dt = 0$ for the first geodesic correction and $\int_{0}^{1}(y' + td)dt = \bar y$. The check solves the problem exactly as a power series in $l$ (Picard iteration with polynomial coefficients in $t$) on two fields in this gauge: a plane wave $\hb^{-1}(a) = a - H(a\cdot k)k$, $k = e_{0} + e_{3}$, $H = \tfrac12(x_{1}^{2} - x_{2}^{2}) + (k\cdot x)x_{1}x_{2}$ (Ricci-flat, Weyl curvature nonzero at $x_{0}$), and a generic polynomial coframe whose linear part is symmetric (so $\Omega(x_{0}) = 0$ by the store's torsion-free solver) with Ricci and Weyl curvature nonzero and $\partial g(x_{0}) \ne 0$; it verifies the two coefficients above (both forms) and the eikonal identity $g^{\mu\nu}\partial_{\mu}\sigma\partial_{\nu}\sigma = 2\sigma$ exactly to order $l^{3}$, with $\partial_{\mu}\sigma = g_{\mu\alpha}(x)\dot\gamma^{\alpha}(1)$.

\emph{Step 3 (the Dirac operator of $\sigma$; proved, computed).} With $g^{\mu} = \hb^{T}(e^{\mu})$ and $\partial_{\mu}\sigma = g_{\mu}\cdot\hb^{-1}(\dot\gamma(1))$ (indices: $g_{\mu\alpha}\dot\gamma^{\alpha}$), $D_{x}\sigma = g^{\mu}\partial_{\mu}\sigma = g^{\mu}\,(g_{\mu}\cdot\hb^{-1}(\dot\gamma(1))) = \hb^{-1}(\dot\gamma(1))$, since $\{g^{\mu}\}$ and $\{g_{\mu}\}$ are reciprocal frames: the Dirac operator of the Synge function is the physical tangent of the geodesic at $x$. With $\dot\gamma = l\dot\eta$: $D_{x}\sigma = l\,\hb^{-1}(ly)(\dot\eta(1)) = l(d + O(l))$. The check verifies the identity against the reciprocal coframe as exact $l$-series.

\emph{Step 4 (Van Vleck determinant and transport rotor; proved, computed).} $\Delta$ is smooth with $\Delta(x, x) = 1$ and, as a biscalar symmetric in its arguments, has no first-order term at the diagonal; its second-order term is $\tfrac16\Ric^{\circ}(x - x', x - x')$ (standard covariant expansion, cited from memory). So $\Delta = 1 + \tfrac{l^{2}}{6}\Ric^{\circ}(d, d) + O(l^{3})$; the check computes $\Delta = -\det(-\partial_{\mu}\partial_{\nu'}\sigma)/\sqrt{\det g(x)\det g(x')}$ from the exact boundary-value problem (an infinitesimal displacement of $y'$) and matches the $l^{2}$ coefficient with the store's Ricci map (gtg.ric\_o) on both fields. For the rotor, $\dot R = -\tfrac{l}{2}\Omega(\dot\eta)(l\eta)R$ gives $R = 1 - \tfrac{l}{2}\Omega(d)(0) + O(l^{2})$, and when $\Omega(0) = 0$, $R = 1 - \tfrac{l^{2}}{2}(\bar y\cdot\nabla)\Omega(d)(0) + O(l^{3})$ (from $\int_{0}^{1}(y' + td)dt = \bar y$); both are verified exactly, with $R\tilde R = 1$ to the computed order. Hence $U = \Delta^{1/2}R = 1 + O(l)$, and $1 + O(l^{2})$ in the gauge $\Omega(x_{0}) = 0$; $U$ is jointly smooth in $(l, y, y')$ and $\partial_{y}U = O(l)$ (in that gauge $O(l^{2})$). In any gauge with $\hb(x_{0}) = 1$, $\partial_{y^{\mu}}R = -\tfrac{l}{2}\Omega(e_{\mu})(0) + O(l^{2})$, so $l^{-1}\partial_{y^{\mu}}R + \tfrac12\Omega(e_{\mu})(ly)R = O(l)$ for every $\mu$: the covariant derivative of the transport rotor vanishes at the diagonal. The check verifies this cancellation exactly with constant bivectors $\Omega(e_{\mu})(0) \ne 0$ added.

\emph{Step 5 (term by term; proved).} The scaled leading term, with $D_{x} = l^{-1}g^{\mu}(ly)\partial_{y^{\mu}} + \tfrac12 g^{\mu}\Omega(e_{\mu})(ly)$ acting from the left:
\[
l^{3}(R_{I_{4}}D_{x} + m_{0})\frac{U}{\sigma} = R_{I_{4}}\Bigl(-\frac{\hb^{-1}(ly)(\dot\eta(1))\,U}{\sigma_{l}^{2}} + \frac{g^{\mu}\partial_{y^{\mu}}U}{\sigma_{l}} + \frac{l}{2}\,\frac{g^{\mu}\Omega(e_{\mu})(ly)\,U}{\sigma_{l}}\Bigr) + \frac{l\,m_{0}U}{\sigma_{l}} .
\]
The first term is $-d/(\tfrac12 Q(d))^{2} + O(l) = -4d/Q(d)^{2} + O(l)$; the second and third are each $O(l)$ when $\Omega(x_{0}) \ne 0$, but by Step 4 their $O(l)$ parts cancel, so their sum $l\,g^{\mu}(l^{-1}\partial_{y^{\mu}}U + \tfrac12\Omega(e_{\mu})(ly)U)/\sigma_{l}$ is $O(l^{2})$ in every gauge with $\hb(x_{0}) = 1$ (each is $O(l^{2})$ when $\Omega(x_{0}) = 0$); the mass term is $O(l)$. Each $O(\cdot)$ is a smooth function of $(l, y, y')$ times the stated power, divided by a power of $\sigma_{l}$. The logarithmic term: $l^{3}D_{x}(V\log\sigma) = l^{2}[(g^{\mu}\partial_{y^{\mu}}V)(\log\sigma_{l} + 2\log l) + V\,\hb^{-1}(\dot\eta(1))/\sigma_{l}] + O(l^{3}\log\sigma_{l})$, with $\partial_{y}V = O(l)$, so it is $O(l^{2})$: $l^{2}$ times a fixed-type distribution ($V(x_{0}, x_{0})\,2d/Q(d)$ at leading order, $-m_{0}^{2}d/Q(d)$ on the flat field) plus $O(l^{3}\log l)$. The $C^{1}$ remainder contributes $l^{3}\,g^{\mu}\partial_{\mu}S$ evaluated at $(ly, ly')$, which is $O(l^{3})$ uniformly on $B\times B$, and its mass term likewise. The check computes these orders (0, 1, 2, 2, 2) as Laurent orders of the exact series on the generic field ($\sigma = l^{2}\sigma_{l}$, $D_{x}\sigma = l\,\hb^{-1}(\dot\eta(1))$, weight $3$), so a wrong power of $l$ in any term fails; for the logarithmic term it uses the flat coefficient $V(x_{0}, x_{0}) = -\tfrac12 m_{0}^{2}$ computed with second-order jets (on a Weyl-curved field $V(x_{0}, x_{0})$ has a curvature part, not computed; only its boundedness is used). It verifies the $l^{0}$ coefficient $-4d/Q(d)^{2}$ of the leading term on both fields, the mass term's $l^{1}$ coefficient $m_{0}/s_{0}$ for $m_{0} = \tfrac12, 3$, and the $l^{1}$ coefficient of the leading term: zero on the plane wave ($\partial g(x_{0}) = 0$), and on the generic field $-(v_{1}/s_{0}^{2} - 2s_{1}d/s_{0}^{3})$ with $s_{0} = \tfrac12 Q(d)$, $s_{1}$ from Step 2 and $v_{1} = (y\cdot\nabla\hb^{-1})(d)|_{0} - \tfrac12\Gamma(0)(d, d)$, $\Gamma(0)(a, b) = (a\cdot\nabla\hb^{-1})(b)|_{0}$ (Step 8).

\emph{Step 6 (the limit is the stored flat massless vacuum; proved).} With $\sigma_{\mathrm{fl}} = \tfrac12 Q(z)$, $z = y - y'$: $\nabla(1/\sigma_{\mathrm{fl}}) = \nabla(2/Q(z)) = -4z/Q(z)^{2}$ and $\nabla(z/Q(z)^{2}) = 4/Q^{2} - 4Q/Q^{3} = 0$ off the cone (verified with exact jets), so the limit of Step 5 is $R_{I_{4}}\bigl((-\tfrac{1}{8\pi^{2}})\cdot(-4z/Q(z)^{2})e_{0}\bigr) = \tfrac{I_{4}}{2\pi^{2}}z e_{0}/Q(z)^{2}$ ($ze_{0}$ is even, so $I_{4}$ on the right and on the left agree), homogeneous of degree $-3$ (weight $\tfrac32$ per leg). Since $-\tfrac{1}{8\pi^{2}}\cdot\tfrac{2}{Q} = -\tfrac{1}{4\pi^{2}Q} = D_{+}$ (Fact \claimref{synthesis.imported-free-field-spin-statistics-and-guido-longo-theorem}, $D_{+}(\xi) = -1/(4\pi^{2}(Q(\xi) - I_{4}0\,\xi_{0}))$), this is $I_{4}\nabla D_{+}(z)e_{0}$, the massless case of the stored kernel $\int\frac{d^{3}p}{(2\pi)^{3}2p_{0}}(p + m_{0})e_{0}e^{-I_{4}p\cdot z}$ of Theorem \claimref{synthesis.bulk-hadamard-states-restrict-to-the-generator-condition-on-flat-null-planes}(ii) ($p\,e^{-I_{4}p\cdot z} = R_{I_{4}}\nabla e^{-I_{4}p\cdot z}$, and $\int\frac{d^{3}p}{(2\pi)^{3}2p_{0}}e^{-I_{4}p\cdot z} = D_{+}(z)$ at $m_{0} = 0$). The other ordering gives $-I_{4}\nabla D_{-}(z)e_{0}$ in the same way. The check verifies the rational coefficients ($-\tfrac18\cdot 2 = -\tfrac14$, $\tfrac14\cdot 2 = \tfrac12$, as coefficients of $\pi^{-2}$) and the homogeneity.

\emph{Step 7 (convergence as distributions; proved, given Fact \claimref{synthesis.imported-hadamard-parametrix-of-the-squared-dirac-operator-leads-with-van-vleck-root-times-spin-transport} (the leading coefficient and the independence of the time function)).} For fixed $l$ the scaled leading distribution is $\lim_{\epsilon\to 0}$ of $1/(\sigma_{l} + I_{4}\epsilon h_{l} + \epsilon^{2})$ times smooth factors, with $h_{l}(y, y') = 2(\tau(ly) - \tau(ly'))/l$ for a time function $\tau$, so that it is the leading Hadamard distribution of the scaled metric $g(ly)$, whose Synge function is $\sigma_{l}$. By Hadamard's lemma $\sigma_{l}(y, y') = d^{\mu}d^{\nu}G_{\mu\nu}(l; y, y')$ with $G$ smooth and symmetric, $G(0) = \tfrac12\eta$; so $\sigma_{l} = \tfrac12 Q(A_{l}d)$ with $A_{l} = (2\eta G)^{1/2}$, the square root near the identity, smooth in $(l, y, y')$, $A_{0} = 1$. The map $\Psi_{l}(y, y') = (y, A_{l}(y, y')d)$ is a diffeomorphism near $B\times B$ for small $l$, and $\Psi_{l} \to \Psi_{0} = (y, y - y')$ in $C^{\infty}$. The scaled leading distribution is the pullback $\Psi_{l}^{*}$ of $1/(\tfrac12 Q(z) + I_{4}0\,h)$ with $h = h_{l}\circ\Psi_{l}^{-1}$ of the sign of $z^{0}$ on the punctured cone; by the independence of the time function this is the pullback of the flat boundary value $1/(\tfrac12 Q(z) - I_{4}0\,\ip{z}{e_{0}})$ (the separation shifted infinitesimally along $-e_{0}$, in scalar form), independent of $l$. A pullback by a diffeomorphism is a change of variables, $\langle\Psi^{*}u, f\rangle = \langle u, (f\circ\Psi^{-1})|\det D\Psi^{-1}|\rangle$, and the test functions on the right converge in $C^{\infty}_{c}$ as $l \to 0$, so $\Psi_{l}^{*}u \to \Psi_{0}^{*}u$ in $\mathcal{D}'$. The measure factor $\sqrt{|\det g|}(ly)$ of the scaled field's pairing is smooth and tends to $1$. Multiplication by smooth factors converging in $C^{\infty}$ and differentiation are continuous, so Step 5 holds in $\mathcal{D}'(B\times B)$, the $O(\cdot)$ terms being smooth functions with the stated powers times distributions of this kind ($1/\sigma_{l}$, $\log\sigma_{l}$ pulled back in the same way). This gives the displayed limit of (iii) for both orderings. For a gauge-invariant quasi-free state the $n$-point functions are sums of products of the two-point functions $\mathcal{G}^{(\pm)}$ (determinants; no anomalous kernels), so they converge too, and the scaling limit is the quasi-free state of the limit kernels, the vacuum of the massless Dirac field on the flat field.

\emph{Step 8 (normal-coordinate dilation; proved, computed).} Let $N(x)$ be the normal coordinates at $x_{0}$, the initial velocity of the geodesic from $x_{0}$ to $x$ with parameter on $[0, 1]$. With $y' = 0$ in Step 2, $\Phi_{l}(y) := l^{-1}N(ly) = \dot\eta(0)$, and $\eta = ty - \tfrac{l}{2}(t^{2} - t)\Gamma(0)(y, y) + O(l^{2})$ gives $\Phi_{l}(y) = y + \tfrac{l}{2}\Gamma(0)(y, y) + O(l^{2})$, where, since $\Omega(0) = 0$ makes $\partial_{\mu}g_{\nu} = \partial_{\nu}g_{\mu}$ at $0$ (torsion-free), $\Gamma(0)(a, b) = (a\cdot\nabla\hb^{-1})(b)|_{0}$ (index reading $\Gamma^{\lambda}{}_{\mu\nu}(0) = e^{\lambda}\cdot\partial_{\mu}g_{\nu}$; checked), smooth in $(l, y)$, tending to the identity with all derivatives. Conservation of $g(\dot\gamma, \dot\gamma)$ along the geodesic gives the Gauss lemma $\sigma(x_{0}, ly) = \tfrac{l^{2}}{2}Q(\Phi_{l}(y))$. The check verifies both on the two fields (the correction vanishes on the plane wave, whose metric has zero first derivatives at $x_{0}$). So test functions dilated by the store's map and by the normal-coordinate map differ by composition with $\Phi_{l} \to \mathrm{id}$, the two scaling limits agree, and (iii) is the Dirac-spinor form of Prop.~5.6 of Sahlmann and Verch, Fact \claimref{synthesis.imported-dirac-hadamard-form-propagates-and-adiabatic-states-are-hadamard}(d) (the statement checked against arXiv math-ph/0008029; their proof in Appendix A.5 not read), the role of their $\Theta$ (Dirac conjugation composed with charge conjugation) being taken by the constant Dirac adjoint $e_{0}$. The physical frame at $x_{0}$ is the background frame because $\hb(x_{0}) = 1$; a constant $\hb(x_{0}) \ne 1$ replaces $y - y'$ by $\hb^{-1}(x_{0})(y - y')$ in the limit.

\emph{Step 9 (what is not given; argued).} The limit kernel restricted to the flat null plane through $x_{0}$ is the null vacuum $W_{0}\otimes\id_{K}$ by Theorem \claimref{synthesis.bulk-hadamard-states-restrict-to-the-generator-condition-on-flat-null-planes}(ii) at $m_{0} = 0$. Clause (b) of (HS) needs more: that the restriction of $\omega$ to the generators of a local Rindler horizon $\mathcal{H}$ through $x_{0}$, scaled along the generators only, tends to it. The steps above do not give this, for two reasons. The restriction to $\mathcal{H}\times\mathcal{H}$ is not a pullback: the wave-front set of the kernel contains $(x, k^{\flat}; x', -k^{\flat})$ for $x, x'$ on one generator, which meets the conormal of $\mathcal{H}\times\mathcal{H}$, so the continuity used in Step 7 (Hormander's pullback theorem, which needs the two to be disjoint) does not apply to the restricted family; on the flat plane the stored theorem defines the restriction by a momentum-space computation that uses translation invariance. And (b) scales the generator parameter with weight $\tfrac12$ and fixed transverse profiles, whereas (iii) scales all four directions with weight $\tfrac32$; passing from one to the other needs the transverse ultralocality of the restricted kernel on $\mathcal{H}$, which is not shown here. Clause (b) is treated by Theorem \claimref{synthesis.field-strength-cancels-from-the-leading-null-singularity-of-hadamard-dirac-states-on-null-hypersurfaces}, which reduces it to a remainder bound on the restricted parametrix, argued; the premise (HS) has since been restated as that bound.
